% !PRINTED 141
\weberpage{141}

Das symbolische Product rechts erhalten wir aber auch, wenn wir
ein Glied der ersten Polare von $\vartheta_x^2 = (b\beta)\,b_x\beta_x$ in die Reihe (VIII)
\S~7 entwickeln. Es ist nämlich:

\weq{2}{(b\beta)\,b_x\beta_y = \{(b\beta)\,b_x\beta_x\}_y + \frac{1}{2}\,(b\beta)^2(xy)\,,}

oder, indem wir nun $x$ durch $a$, $y$ durch $\alpha$ ersetzen und die Symbole
$\vartheta_x^2$, $A_{f\varphi}$ für $(b\beta)\,b_x\beta_x$, resp. $(b\beta)^2$ einführen und mit $(a\alpha)$
multipliciren:

\weq{3}{(b\beta)(ba)(\beta\alpha)(a\alpha) = (\vartheta a)(\vartheta\alpha)(a\alpha) + \frac{1}{2}\,A_{f\varphi}\cdot(a\alpha)^2\,.}

Hierin tritt bereits das symbolische Product der Gleichung (1)
auf, und da $(a\alpha)^2 = A_{f\varphi}$, so handelt es sich nur mehr darum, auch
die linke Seite dieser Gleichung (3) auf bekannte Formen zurückzuführen. Es ist aber, wenn wir darin $a$ mit $b$ vertauschen, und die
halbe Summe der so entstehenden Producte nehmen:

\[
(b\beta)(ba)(\beta\alpha)(a\alpha) = \frac{1}{2}\,(ab)(\alpha\beta)\{(a\alpha)(b\beta) - (b\alpha)(a\beta)\}
\]
oder

\weq{4}{= \frac{1}{2}\,(ab)^2(\alpha\beta)^2 = \frac{1}{2}\,A_{ff}\cdot A_{\varphi\varphi}\,,}

wie der Identitätssatz unmittelbar liefert. Demnach erhalten wir aus
Gleichung (3) und (4)

\weq{5}{\frac{1}{2}\,\{A_{ff}\cdot A_{\varphi\varphi} - A_{f\varphi}^2\} = (a\alpha)(\vartheta\alpha)(\vartheta a)\,,}

oder durch Vergleichung mit (1)

\weq{6}{(\vartheta\vartheta_1)^2 = \frac{1}{2}\,\{A_{ff}A_{\varphi\varphi} - A_{f\varphi}^2\}\,,}

d. h.: \glqq Jedes symbolische Product mit dem Factor $(\vartheta\vartheta_1)^2$ ist
reducibel.\grqq{}
Dies gilt aber auch für symbolische Producte mit den Factoren
$(a\vartheta)^2$ und $(\alpha\vartheta)^2$. Denn es ist:

\weq{7}{
\left.
\begin{aligned}
0 &= A_{f\vartheta} = (a\vartheta)^2 = (b\alpha)(ab)(a\alpha) \\
0 &= A_{\varphi\vartheta} = (\alpha\vartheta)^2 = (a\beta)(\alpha a)(\alpha\beta)
\end{aligned}
\right\}\,,
}

wie man unmittelbar erkennt, da die symbolischen Producte rechts
durch Vertauschung gleichberechtigter Symbole nur ihr Vorzeichen
ändern. Wir haben somit den Satz bewiesen:

\begin{quote}
\glqq Alle symbolischen Producte, welche einen der sechs möglichen
Klammerfactoren \gesperrt{quadratisch} enthalten, sind reducibel.\grqq{}
\end{quote}

Es erübrigt also nur noch, dieselbe Eigenschaft für symbolische
Producte nachzuweisen, welche diese Factoren linear besitzen.

126. \textit{Beweis, dass alle Producte, welche einen der sechs Factoren (A)
in Nr. 123 linear enthalten, reducibel sind.}


% ===== pdf idx 156, printed 142.0 =====
% !PRINTED 142
\weberpage{142}

Die Factoren $(ab)$ und $(\alpha\beta)$ haben wir bereits in Nr. 119 als Reducenten
erkannt; genau dieselbe Eigenschaft besitzt $(\vartheta\vartheta_1)$ aus gleichen Gründen.
Es bleiben also nur noch Producte mit den linearen Klammerfactoren
$(a\alpha)$, $(a\vartheta)$, $(\alpha\vartheta)$ zu untersuchen. Ein Product, das den Factor $(a\vartheta)$
oder $(\alpha\vartheta)$ besitzt, muss auch noch die Factoren $a_\xi\vartheta_\eta$ resp.
$\alpha_\xi\vartheta_\eta$ enthalten, wobei die Grössen $\xi$, $\eta$ für irgend welche andere
Symbole oder Variable stehen mögen. Wir brauchen also nur zu zeigen, dass

\[
(a\vartheta)\,a_\xi\,\vartheta_\eta \quad \text{und} \quad (\alpha\vartheta)\,\alpha_\xi\,\vartheta_\eta
\]

reducibel sind. Dies erkennen wir aber auf folgende Weise. Da $(a\vartheta)$
und $(\alpha\vartheta)$ die Klammerfactoren der Functionaldeterminanten $(f,\ \vartheta)$
und $(\varphi,\ \vartheta)$ sind, so sind sie sicher Reducenten, wenn diese Formen
und die aus ihnen durch Faltung entstehenden reducibel sind. Es
sind aber $(f,\ \vartheta)$ und $(\varphi,\ \vartheta)$ Functionaldeterminanten von Functionaldeterminanten, und als solche nach den früher gegebenen allgemeinen
Theorien stets durch einfachere Formen darstellbar. In der That
liefert Nr. 52, (I) und (II) direct:

\[
(f,\ \vartheta) = \frac{1}{2}\,\{A_{f\varphi}\cdot f - A_{ff}\cdot\varphi\}
\]

\[
(\vartheta,\ \varphi) = \frac{1}{2}\,\{A_{f\varphi}\cdot\varphi - A_{\varphi\varphi}\cdot f\}\,.
\]

Die aus $(f,\ \vartheta)$ und $(\vartheta,\ \varphi)$ durch Faltung entstehenden Formen
$A_{f\vartheta}$ und $A_{\varphi\vartheta}$ sind aber gleichfalls reducibel, da sie verschwinden;
demnach sind $(a\vartheta)$ und $(\alpha\vartheta)$ in der That Reducenten.

Man kann die Reducibilität der Formen $(a\vartheta)\,a_\xi\vartheta_\eta$ und $(\alpha\vartheta)\,\alpha_\xi\vartheta_\eta$
auch in anderer Weise einfach zeigen. Aus den Formeln Nr. 51, (II)
erhält man nämlich direct:

\[
(a\vartheta)\,\vartheta_\eta\,a_\xi = (f,\ \vartheta)_\eta = \frac{1}{2}\,\{A_{f\varphi}\cdot f_\xi\cdot f_\eta - A_{ff}\ \cdot\varphi_\xi\cdot\varphi_\eta\}
\]

\[
-(\alpha\vartheta)\,\vartheta_\eta\,\alpha_\xi = (\vartheta,\ \varphi)_\eta = \frac{1}{2}\,\{A_{f\varphi}\cdot\varphi_\xi\cdot\varphi_\eta - A_{\varphi\varphi}\cdot f_\xi\cdot f_\eta\}\,.
\]

Die symbolischen Producte rechts enthalten neben den bekannten
drei Invarianten des Systemes nur mehr ein Symbol $a$, respective
$\alpha$, also \gesperrt{weniger} Symbole wie die Ausdrücke links, wodurch deren
Reducibilität nachgewiesen ist.

Es erübrigt noch, die Reducibilität von Producten mit dem Klammerfactor $(a\alpha)$, oder, was dasselbe ist, mit dem Theilproduct $(a\alpha)\,a_\xi\alpha_\eta$
nachzuweisen. Sie ergiebt sich direct aus Nr. 125, (2), wonach:

\[
(a\alpha)\,a_\xi\,\alpha_\eta = \vartheta_\xi\,\vartheta_\eta + \frac{1}{2}\,A_{f\varphi}\,(\xi\eta)\,.
\]


% ===== pdf idx 157, printed 143.0 =====
% !PRINTED 143
\weberpage{143}

Die rechte Seite enthält nur reducible Glieder, das zweite wegen
des Factors $A_{f\varphi}$, das erste, da wir für jeden Klammerfactor mit dem
Symbol $\vartheta$ bereits die Eigenschaft nachgewiesen haben, ein Product,
dem er angehört, reducibel zu machen.

Wir sind also damit zu dem Resultate gelangt, dass jedes symbolische Product, welches irgend einen der sechs möglichen Klammerfactoren linear enthält, auf die sechs bekannten Formen reducirt
werden kann.

\begin{quote}
\glqq Demnach sind alle sechs Klammerfactoren

\[
(ab),\ (\alpha\beta),\ (a\alpha);\quad (a\vartheta),\ (\alpha\vartheta),\ (\vartheta\vartheta)
\]

Reducenten, und es lassen sich alle Co- und Invarianten der
Formen $f$ und $\varphi$ ganz und rational durch die Formen

\[
f,\ \varphi,\ \vartheta,\ A_{ff},\ A_{\varphi\varphi},\ A_{f\varphi}
\]

darstellen.\grqq{}
\end{quote}

127. \textit{Zusammenfassung der gewonnenen Relationen.} Ich will nur
noch tabellarisch auf die wichtigeren Relationen, die bei dieser Untersuchung sich ergeben haben, aufmerksam machen, indem ich sie im
Folgenden zusammenstelle. Ist $f = a_x^2$, $\varphi = \alpha_x^2$, $\vartheta = (a\alpha)\,a_x\,\alpha_x$, so
hat man:

\[
\begin{array}{rll}
(1) & 2A_{\vartheta\vartheta} = A_{ff}A_{\varphi\varphi} - A_{f\varphi}^2 & \text{(Nr. 125)} \\[4pt]
(2) & A_{f\vartheta} = 0,\quad A_{\varphi\vartheta} = 0 & \text{(Nr. 125)} \\[4pt]
(3) & 2(\vartheta,\ f) = A_{ff}\cdot\varphi - A_{f\varphi}\cdot f & \text{(Nr. 126)} \\[4pt]
    & 2(\varphi,\ \vartheta) = A_{\varphi\varphi}\cdot f - A_{f\varphi}\cdot\varphi. &
\end{array}
\]

Hiezu tritt:

\[
(4)\quad -2\vartheta^2 = A_{ff}\varphi^2 - 2A_{f\varphi}f\varphi + A_{\varphi\varphi}f^2\,,
\]

eine Relation (4), welche zwar hier nicht neuerdings abgeleitet wurde,
die sich aber unmittelbar aus der allgemeineren Nr. 48, (I) ergiebt.

128. \textit{Erste Anwendung.} Ist $\vartheta$ identisch Null, so sind $f$ und
$\varphi$ einander proportional. Wir zeigen zuerst, dass $\vartheta = 0$ die hinreichende Bedingung ist, damit die Gleichung besteht

\[
f = \varrho \cdot \varphi\,,
\]

wo $\varrho$ ein Proportionalitätsfactor. Es ist nämlich

\[
2\vartheta_x\vartheta_y = (a\alpha)\{a_x\alpha_y + \alpha_x a_y\}\,,
\]

und weil nach dem Productsatze:

\[
(xy) = \frac{1}{(a\alpha)}\{a_x\alpha_y - \alpha_x a_y\}\,,
\]

so folgt durch Multiplication beider Gleichungen

\[
2(xy)\vartheta_x\vartheta_y = a_x^2\alpha_x^2 - a_y^2\alpha_x^2\,.
\]


% ===== pdf idx 158, printed 144.0 =====
% !PRINTED 144
\weberpage{144}

Wenn nun $\vartheta = 0$, dann ist auch $\vartheta_y = 0$, also wird in diesem Falle:

\[
a_x^2\,\alpha_y^2 = a_y^2\,\alpha_x^2\,,
\]

oder

\[
a_x^2 = \frac{a_y^2}{\alpha_y^2}\cdot\alpha_x^2 = \varrho\cdot\alpha_x^2\,.
\]

Die Bedingung $\vartheta = 0$ ist aber nicht nur hinreichend, sondern
auch nothwendig. Denn wenn

\[
\alpha_x^2 = \frac{1}{\varrho}\,a_x^2 = \frac{1}{\varrho}\,b_x^2\,,
\]

dann ist

\[
\vartheta = (a_x^2,\ \alpha_x^2) = \frac{1}{\varrho}\,(a_x^2,\ b_x^2) = \frac{1}{\varrho}\,(ab)\,a_x\,b_x = 0\,.
\]

129. \textit{Zweite Anwendung.} Ist $A_{\vartheta\vartheta}$ identisch null, so besitzen $f$ und
$\varphi$ einen gemeinsamen linearen Factor. Diese Bedingung ist nothwendig; denn wenn $f$ und $\varphi$ einen gemeinsamen linearen Factor besitzen, also etwa

\[
f = p_x\,q_x
\]

\[
\varphi = p_x\,r_x
\]

ist, dann wird $\vartheta$ ein reines Quadrat. Denn man hat nach Nr. 38, (I)

\[
\vartheta = (p_x\,q_x,\ p_x\,r_x) = \frac{1}{4}\,\{(pp)\,q_x\,r_x + (pr)\,q_x\,p_x + (qp)\,p_x\,r_x + (qr)\,p_x^2\}\,.
\]

Da aber $(pp) = 0$ und $(pr)\,q_x + (qp)\,r_x = (qr)\,p_x$, so wird

\[
\vartheta = \frac{1}{4}\,\{(qr)\,p_x^2 + (qr)\,p_x^2\}\,,
\]

oder

\[
\vartheta = \frac{1}{2}\,(qr)\cdot p_x^2\,.
\]

Die Functionaldeterminante $\vartheta$ ist also das Quadrat des gemeinsamen
linearen Factors $p_x$ und demnach $A_{\vartheta\vartheta}$ als Discriminante identisch null.

Die Bedingung $A_{\vartheta\vartheta} = 0$ ist aber auch hinreichend. Denn wenn

\[
A_{\vartheta\vartheta} = 0\,,
\]

dann ist

\[
\vartheta = m_x^2\,,
\]

wo $m_x$ eine lineare Form ist. Die zweite Ueberschiebung von $f = a_x^2$
über $\vartheta = m_x^2$ ist sodann

\[
A_{f\vartheta} = (am)^2 = a_m^2\,.
\]

Wir wissen aber, dass $A_{f\vartheta}$ wie $A_{\varphi\vartheta}$ identisch verschwinden; es
ist also sowohl

\[
a_m^2 = 0\,,
\]

als auch

\[
\alpha_m^2 = 0\,,
\]
%CONT>>


% ===== pdf idx 159, printed 145.0 =====
% !PRINTED 145
\weberpage{145}
%<<CONT
d. h. die Formen $a_x^2$ und $\alpha_x^2$ werden identisch null für $x=m$, oder was das
Nämliche ist, sie besitzen einen gemeinschaftlichen linearen Factor $m_x$.

Da also $A_{\vartheta\vartheta}=0$ die nothwendige und hinreichende Bedingung ist,
dass $f$ und $\varphi$ einen gemeinsamen Factor besitzen, so muss $A_{\vartheta\vartheta}$
der Resultante $R_{f,\varphi}$ beider Formen proportional sein. Die
Proportionalitätsconstante kann nur ein Zahlenfactor sein, da sowohl $A_{\vartheta\vartheta}$
als auch $R_{f,\varphi}$ vom vierten Grade in den Coefficienten beider Formen ist. Man
findet durch ein beliebiges Beispiel leicht:

\[
R_{f,\varphi} = -2A_{\vartheta\vartheta} = A_{f\varphi}^2 - A_{ff}A_{\varphi\varphi}.
\]

\textit{Anmerkung.} Die Identität:

\[
2\vartheta = (qr)\cdot p_x^2
\]
lehrt zugleich, wie man den gemeinschaftlichen Factor ermitteln kann. Denn vergleicht
man in der ersten Polare

\[
2\vartheta_x\vartheta_y = (qr)p_xp_y
\]
die Coefficienten von $y_1$ und $y_2$, so erhält man:

\[
\begin{gathered}
2\vartheta_x\vartheta_1 = (qr)p_xp_1\\
2\vartheta_x\vartheta_2 = (qr)p_xp_2,
\end{gathered}
\]
und somit durch Division:

\weq{1}{\vartheta_x\vartheta_1 : \vartheta_x\vartheta_2 = p_1:p_2}

für jeden Werth von $x$. Die beiden ersten partiellen Differentialquotienten von
$\vartheta$ sind sonach den Coefficienten des gemeinschaftlichen linearen Factors von
$f$ und $\varphi$ proportional; dies gilt insbesondere beispielsweise für $x_1=1,
x_2=0$, für welchen Fall Gleichung (1) in

\weq{2}{\bar\vartheta_0 : \bar\vartheta_1 = p_1 : p_2}

übergeht. Es ist aber nach Nr. 47, (II)

\[
\vartheta = \begin{vmatrix} \bar a_0 & \bar a_1 & \bar a_2 \\ \bar\alpha_0 & \alpha_1 & \alpha_2 \\ x_1^2 & -x_1x_2 & x_2^2 \end{vmatrix} = \bar\vartheta_0 x_1^2 + 2\bar\vartheta_1 x_1x_2 + \bar\vartheta_2 x_2^2,
\]
wobei

\[
\bar\vartheta_0 = (\bar a_1\bar\alpha_2), \quad \bar\vartheta_1 = \frac{1}{2}(\bar a_0\bar\alpha_2), \quad \bar\vartheta_2=(\bar a_0\bar\alpha_1).
\]

Durch Substitution dieser Werthe in (2) erhalten wir:

\weq{3}{\frac{p_1}{p_2} = \frac{2(\bar a_1\bar\alpha_2)}{(\bar a_0\bar\alpha_2)}}

130. \textit{Dritte Anwendung. Kanonische Form für zwei simultane
quadratische Formen.} Wir stellen uns die Aufgabe, zwei solche lineare
%CONT>>


% ===== pdf idx 160, printed 146.0 =====
% !PRINTED 146
\weberpage{146}
%<<CONT
Formen $r_x$ und $s_x$ zu ermitteln, dass durch dieselbe Transformation $f$ und
$\varphi$ gleichzeitig die Gestalt erhalten:

\weq{I}{\left.\begin{aligned} f &= \lambda_1 r_x^2 + \lambda_2 s_x^2\\ \varphi &= \mu_1 r_x^2 + \mu_2 s_x^2\end{aligned}\right\}}

Die erste Untersuchung zur Lösung solcher Aufgaben ist immer: Wie gestaltet sich
das simultane Formensystem, wenn wir $f$ und $\varphi$ bereits in dieser kanonischen
Form zu Grunde legen? Es ergiebt sich sofort, dass $\vartheta$ das Product der beiden
linearen Formen $r_x$ und $s_x$ ist, bis auf einen constanten Factor. Denn man erhält:

\[
\begin{gathered}
\vartheta = (f,\varphi) = (\lambda_1 r_x^2 + \lambda_2 s_x^3,\ \mu_1 r_x^2 + \mu_2 s_x^2) % [sic]
\\
= \lambda_1\mu_1(r_x^2,r_x^2) + \lambda_1\mu_2(r_x^2,s_x^2) + \lambda_2\mu_1(s_x^2,r_x^2) + \lambda_2\mu_2(s_x^2,s_x^2)
\\
= \lambda_1\mu_2(rs)r_xs_x - \lambda_2\mu_1(rs)r_xs_x
\\
= (\lambda\mu)(rs)r_xs_x.
\end{gathered}
\]

Wir lernen daraus: Die Factoren $r_x$ und $s_x$, durch welche sich $f$ und $\varphi$
gleichzeitig in die kanonische Form (I) bringen lassen, sind den linearen Factoren
von $\vartheta$ proportional. Damit also überhaupt diese Darstellung möglich ist,
muss $\vartheta$ wirklich zwei verschiedene lineare Factoren besitzen, d. h. $f$ und
$\varphi$ dürfen keinen Factor gemeinsam haben.

Die Einführung dieser beiden linearen Factoren von $\vartheta$ ist aber auch
hinreichend, um $f$ und $\varphi$ gleichzeitig als Summe zweier Quadrate darzustellen.
Denn sei

\[
\vartheta = r_x \cdot s_x,
\]
dann ist nach Nr. 126 (2):

\weq{1}{\left.\begin{aligned} A_{f\vartheta} &= (as)(ar) = 0\\ A_{\varphi\vartheta} &= (\alpha s)(\alpha r) = 0\end{aligned}\right\}}

Nun hat man gemäss dem Identitätssatze:

\[
\begin{gathered}
a_x(rs) = r_x(as) - s_x(ar)\\
\alpha_x(rs) = r_x(\alpha s) - s_x(\alpha r).
\end{gathered}
\]

Also, indem man beide Identitäten quadrirt, und berücksichtigt, dass wegen der
Relationen (1) die mittleren Glieder rechts verschwinden:

\weq{II}{\left.\begin{aligned} f\cdot(rs)^2 &= r_x^2(as)^2 + s_x^2(ar)^2\\ \varphi\cdot(rs)^2 &= r_x^2(\alpha s)^2 + s_x^2(\alpha r)^2\end{aligned}\right\}}

wodurch in der That $f$ und $\varphi$ als Summen zweier Quadrate dargestellt sind.

Um die Transformation von $f$ und $\varphi$ wirklich auszuführen, wird man sonach
entweder die Form $\vartheta$ in ihre Linearfactoren spalten oder
%CONT>>


% ===== pdf idx 161, printed 147.0 =====
% !PRINTED 147
\weberpage{147}
%<<CONT
einfacher in folgender Weise verfahren. Man multiplicirt die Gleichungen (II) mit
1 resp. $\lambda$ und addirt dieselben; dann erhält man:

\[
f + \lambda\varphi = r_x^2\left\{\frac{(as)^2+\lambda(\alpha s)^2}{(rs)^2}\right\} + s_x^2\left\{\frac{(ar)^2+\lambda(\alpha r)^2}{(rs)^2}\right\}.
\]

Nun lässt sich $\lambda$ so bestimmen, dass der Coefficient von $r_x^2$ oder von
$s_x^2$ verschwindet; in diesem Falle geht dann $f+\lambda\varphi$ gerade in das
Quadrat der linearen Form $s_x$ resp. $r_x$ (bis auf einen constanten Factor) über.
Damit aber $f+\lambda\varphi$ ein reines Quadrat sei, muss die Discriminante:

\[
(f+\lambda\varphi,\ f+\lambda\varphi)^2 = (f,f)^2 + 2\lambda(f,\varphi)^2 + \lambda^2(\varphi,\varphi)^2
\]
identisch verschwinden. Dies liefert eine quadratische Gleichung für $\lambda$,
nämlich:

\[
A_{ff} + 2\lambda A_{f\varphi} + \lambda^2 A_{\varphi\varphi} = 0.
\]

Sind dann $\lambda_1$ und $\lambda_2$ die Wurzeln dieser Gleichung, so ist:

\[
\begin{gathered}
f + \lambda_1\varphi = m_x^2\\
f + \lambda_2\varphi = n_x^2,
\end{gathered}
\]
und somit:

\[
\begin{gathered}
f = \frac{\lambda_2}{\lambda_2-\lambda_1}m_x^2 - \frac{\lambda_1}{\lambda_2-\lambda_1}n_x^2\\
\varphi = \frac{1}{\lambda_1-\lambda_2}m_x^2 - \frac{1}{\lambda_1-\lambda_2}n_x^2.
\end{gathered}
\]

\begin{center}
\textbf{\S\ 12. Die Form zweiten Grades: System dreier und mehr simultaner Formen.}
\end{center}

131. \textit{Das simultane System dreier quadratischer Formen.} Das volle System
dreier quadratischer Formen $f,\varphi,\psi$ besteht aus folgenden Bildungen.

1) Die Originalformen selbst:

\[
\begin{gathered}
f = a_x^2 = b_x^2 = \dots\\
\varphi = \alpha_x^2 = \beta_x^2 = \dots\\
\psi = r_x^2 = s_x^2 = \dots;
\end{gathered}
\]

2) ihre sechs zweiten Ueberschiebungen:

\[
\begin{gathered}
(f,f)^2 = A_{ff}, \quad (\varphi,\varphi)^2 = A_{\varphi\varphi}\\
(f,\varphi)^2 = A_{f\varphi}, \quad (\varphi,\psi)^2 = A_{\varphi\psi}\\
(f,\psi)^2 = A_{f\psi}, \quad (\psi,\psi)^2 = A_{\psi\psi};
\end{gathered}
\]

3) ihre drei Functionaldeterminanten (schiefe Covarianten):

\[
\begin{gathered}
(\varphi,\psi) = \vartheta_1 = \vartheta_{1x}^2\\
(\psi,f) = \vartheta_2 = \vartheta_{2x}^2\\
(f,\varphi) = \vartheta_3 = \vartheta_{3x}^2;
\end{gathered}
\]


% ===== pdf idx 162, printed 148.0 =====
% !PRINTED 148
\weberpage{148}
4) ihre Combinante (die einzige schiefe Invariante):

\[
(a\alpha)(\alpha r)(ar) = \begin{vmatrix} \bar a_0 & \bar a_1 & \bar a_2 \\ \bar\alpha_0 & \bar\alpha_1 & \bar\alpha_2 \\ \bar r_0 & \bar r_1 & \bar r_2 \end{vmatrix} = R_{123},
\]
welch letztere wir bereits in § 1 kennen gelernt haben. Ausser diesen dreizehn
Formen existirt keine, die sich nicht als ganze rationale Function derselben
darstellen liesse.

132. \textit{Alle Producte mit quadratischen Klammerfactoren sind auf diese
dreizehn Formen reducibel.} Irgend ein symbolisches Product nämlich, das eine neue
Covariante der drei Formen repräsentiren soll, kann nur folgende Klammerfactoren

\weq{A}{\begin{gathered}
1)\quad (ab)\quad (\alpha\beta)\quad (rs)\quad (a\alpha)\quad (ar)\quad (\alpha r)\\
(a\vartheta_3)\quad (a\vartheta_2)\quad (\alpha\vartheta_1)\quad (\alpha\vartheta_3)\quad (r\vartheta_2)\quad (r\vartheta_1)\\
(\vartheta_1\vartheta_1)\quad (\vartheta_2\vartheta_2)\quad (\vartheta_3\vartheta_3);\\
2)\quad (a\vartheta_1)\quad (\alpha\vartheta_2)\quad (r\vartheta_3)\quad (\vartheta_1\vartheta_2)\quad (\vartheta_1\vartheta_3)\quad (\vartheta_2\vartheta_3)
\end{gathered}}

besitzen und zwar im höchsten Falle in der zweiten Potenz. Was die fünfzehn
Factoren der ersten Gruppe betrifft, so sahen wir bereits bei Aufstellung des
simultanen Systemes zweier quadratischer Formen, dass jedes Product, in welchem
sie quadratisch oder linear auftreten, auf die Formen des Systemes reducibel ist.
Unsere Aufgabe ist also nur noch nachzuweisen, dass auch die Klammerfactoren der
zweiten Gruppe Reducenten sind. Wir betrachten wiederum zuerst Producte, deren
Klammerfactoren den Exponenten 2 haben. Die Factoren

\[
(a\vartheta_1)^2,\ (\alpha\vartheta_2)^2,\ (r\vartheta_3)^2,\ (\vartheta_1\vartheta_2)^2,\ (\vartheta_1\vartheta_3)^2,\ (\vartheta_2\vartheta_3)^2
\]
gehören zunächst folgenden Ueberschiebungen an:

\weq{1}{(f,(\varphi,\psi))^2,\qquad (\varphi,(\psi,f))^2,\qquad (\psi,(f,\varphi))^2}

\weq{2}{((\varphi,\psi),(\psi,f))^2,\quad ((\varphi,\psi),(f,\varphi))^2,\quad ((\psi,f),(f,\varphi))^2.}

Es ist aber:

\weq{I}{\begin{aligned} (f,(\varphi,\psi))^2 &= (a_x^2,(\alpha r)\alpha_xr_x)^2 = (a\alpha)(\alpha r)(ar) = R_{123}\\ &= (\varphi,(\psi,f))^2 = (\psi,(f,\varphi))^2,\end{aligned}}

und folglich führen die Ueberschiebungen der Gruppe (1) auf keine neue Form, d. h.
Producte mit den Factoren $(a\vartheta_1)^2$, $(\alpha\vartheta_2)^2$,
$(r\vartheta_3)^2$ sind reducibel.

Was die Ueberschiebungen der Gruppe (2) betrifft, so können wir deren
Reducibilität auf folgende Art beweisen. Ich will zuerst den Weg der directen
Rechnung einschlagen. Es ist, wenn $f=a_x^2$, $\varphi=\alpha_x^2$, $\psi=r_x^2$,
$\chi=\mu_x^2$ irgend vier quadratische Formen darstellen:
%CONT>>


% ===== pdf idx 163, printed 149.0 =====
% !PRINTED 149
\weberpage{149}

\[
\begin{aligned}
U &= ((f,\varphi),(\psi,\chi))^2 = (a\alpha)(r\mu)(a_x\alpha_x,\ r_x\mu_x)^2\\
&= \frac{1}{2}(a\alpha)(r\mu)\{(ar)(\alpha\mu)+(a\mu)(\alpha r)\}\\
&= \frac{1}{2}(G_1+G_2).
\end{aligned}
\]
Die Differenz der beiden Glieder rechts ist:
\[
G_1-G_2=(a\alpha)(r\mu)\{(ar)(\alpha\mu)-(a\mu)(\alpha r)\},
\]
oder weil:
\[
(ar)(\alpha\mu)-(a\mu)(\alpha r)=(a\alpha)(r\mu)
\]
(Identit\"atssatz), so erh\"alt man:
\[
G_1-G_2=(a\alpha)^2(r\mu)^2=A_{f\varphi}A_{\psi\chi}.
\]
Addiren wir die beiden Gleichungen:
\[
\begin{aligned}
U &= \frac{1}{2}(G_1+G_2)\\
\frac{1}{2}A_{f\varphi}A_{\psi\chi} &= \frac{1}{2}(G_1-G_2),
\end{aligned}
\]
so erh\"alt man:
\[
U+\frac{1}{2}A_{f\varphi}A_{\psi\chi}=(a\alpha)(r\mu)(ar)(\alpha\mu),
\]
oder, wenn wir auf das symbolische Product rechts den Productsatz anwenden:
\[
U+\frac{1}{2}A_{f\varphi}A_{\psi\chi}=\frac{1}{2}\{A_{f\varphi}A_{\psi\chi}+A_{f\psi}A_{\varphi\chi}-A_{f\chi}A_{\varphi\psi}\},
\]
also
\weq{II}{U=((f,\varphi),(\psi,\chi))^2=\frac{1}{2}\{A_{f\psi}A_{\varphi\chi}-A_{f\chi}A_{\varphi\psi}\}.}
F\"ur $\chi=f$ erh\"alt man hieraus
\[
((f,\varphi),(\psi,f))^2=\frac{1}{2}(A_{f\psi}A_{f\varphi}-A_{ff}A_{\varphi\psi})=A_{\vartheta_3\vartheta_2}.
\]
Man kann das gleiche Resultat rascher durch den \gesperrt{Aronhold}'schen Process erlangen. F\"uhrt man n\"amlich in die Identit\"at
\[
\begin{aligned}
2(\vartheta_3,\vartheta_3)^2 &= A_{ff}A_{\varphi\varphi}-A_{f\varphi}^2\\
&= (ab)^2(\alpha\beta)^2-(a\alpha)^2(b\beta)^2
\end{aligned}
\]
durch einmaliges Deltairen statt des Symboles $\alpha$ das Symbol $r$ ein, so kommt, weil
\[
\begin{aligned}
\delta(\vartheta_3) &= \delta(f,\varphi)=(f,\psi)=-\vartheta_2\\
-4(\vartheta_3,\vartheta_2)^2 &= 2(ab)^2(\alpha r)^2-2(a\alpha)^2(br)^2,
\end{aligned}
\]
oder
\weq{III}{A_{\vartheta_3\vartheta_2}=-\frac{1}{2}\{A_{ff}A_{\varphi\psi}-A_{f\varphi}A_{f\psi}\}=(\vartheta_3\vartheta_2)^2.}
Analoge Relationen erh\"alt man f\"ur $A_{\vartheta_3\vartheta_1}$ und $A_{\vartheta_1\vartheta_2}$. Folglich sind Producte, welche die Factoren $(\vartheta_1\vartheta_2)^2$, $(\vartheta_1\vartheta_3)^2$, $(\vartheta_2\vartheta_3)^2$ besitzen, auf Formen des Systems reducibel.


% ===== pdf idx 164, printed 150.0 =====
% !PRINTED 150
\weberpage{150}

133. \gesperrt{Producte, welche die Klammerfactoren (A) linear enthalten, sind reducibel.} Es bleiben also nur mehr Producte zu untersuchen, in denen kein Klammerfactor in h\"oherer als in der ersten Potenz auftritt. Auch hier brauchen wir die Reducibilit\"at nur mehr f\"ur Producte nachzuweisen, in welchen die letzten sechs Klammerfactoren
\[
(a\vartheta_1)\quad(\alpha\vartheta_2)\quad(r\vartheta_3)\quad(\vartheta_1\vartheta_2)\quad(\vartheta_1\vartheta_3)\quad(\vartheta_2\vartheta_3)
\]
linear auftreten. Der Beweis gestaltet sich genau wie bei zwei simultanen quadratischen Formen. Wir haben nur zu zeigen, dass die Producte
\[
(a\vartheta_1)a_x\vartheta_{1y},\ \ldots,\ (\vartheta_1\vartheta_2)\vartheta_{1x}\vartheta_{2y},\ \ldots
\]
reducibel sind. Die hiezu n\"othigen Formeln gehen direct aus der Relation (2) Nr. 125
\[
(f,\varphi)_y=(a\alpha)a_x\alpha_y-\frac{1}{2}A_{f\varphi}(xy)
\]
hervor. Ersetzt man hierin n\"amlich die quadratische Form $\varphi$ durch die quadratische Form $\vartheta_1$, so kommt:
\weq{I}{(f,\vartheta_1)_y=(a\vartheta_1)a_x\vartheta_{1y}-\frac{1}{2}A_{f\vartheta_1}(xy);}
und substituirt man hierin noch $\vartheta_2$ f\"ur $f$, so kommt auch:
\weq{II}{(\vartheta_2,\vartheta_1)_y=(\vartheta_2\vartheta_1)\vartheta_{2x}\vartheta_{1y}-\frac{1}{2}A_{\vartheta_2\vartheta_1}(xy).}
Beide Relationen lehren, dass die Producte
\[
(a\vartheta_1)a_x\vartheta_{1y}\quad\text{und}\quad(\vartheta_2\vartheta_1)\vartheta_{2x}\vartheta_{1y}
\]
reducibel sind; denn die Formen $A_{f\vartheta_1}$, $A_{\vartheta_2\vartheta_1}$ sind reducibel gem\"ass den Entwicklungen in Nr. 132, und die Formen $(f,\vartheta_1)$ und $(\vartheta_2,\vartheta_1)$ sind als Functionaldeterminanten von Functionaldeterminanten reducibel, und zwar hat man, wie aus der allgemeinen Formel (III) Nr. 53 direct hervorgeht:
\weq{III}{(f,\vartheta_1)=(f,(\varphi,\psi))=-\frac{1}{2}\{A_{f\varphi}\psi-A_{f\psi}\varphi\},}
und ebenso:
\[
(\vartheta_2,\vartheta_1)=((\psi,f),\vartheta_1)=-\frac{1}{2}\{A_{f\vartheta_1}\psi-A_{\psi\vartheta_1}f\},
\]
oder weil:
\[
A_{\psi\vartheta_1}=(r\vartheta_1)^2=(\psi,(\varphi,\psi))^2=0
\]
und
\[
A_{f\vartheta_1}=(a\vartheta_1)^2=R_{123}:
\]
so wird:
\weq{IV}{(\vartheta_2,\vartheta_1)=-\frac{1}{2}R_{123}\psi.}
Die Relationen (I) und (II) gehen also unter Benutzung dieser Werthe (III) und (IV) \"uber in:
%CONT>>


% ===== pdf idx 165, printed 151.0 =====
% !PRINTED 151
\weberpage{151}
%<<CONT
\weq{V}{
\begin{aligned}
(a\vartheta_1)a_x\vartheta_{1y} &= \frac{1}{2}\{A_{f\psi}\varphi_y\varphi_x-A_{f\varphi}\psi_y\psi_x+R_{123}(xy)\}\\
(\vartheta_1\vartheta_2)\vartheta_{2x}\vartheta_{1y} &= \frac{1}{2}\{R_{123}\psi_y\psi_x+\frac{1}{2}(xy)(A_{\psi\psi}A_{f\varphi}-A_{\psi\varphi}A_{\psi f})\}
\end{aligned}
}
Somit sind alle Producte, welche irgend einen der vorhandenen Klammerfactoren \gesperrt{linear} enthalten, reducibel. Durch Faltung derselben entstehen Producte mit quadratischen Klammerfactoren, die gleichfalls auf die Formen des Systemes zur\"uckgef\"uhrt werden k\"onnen. Demnach sind alle Klammerfactoren Reducenten und das simultane System dreier quadratischer Formen besteht aus keinen weiteren als den Eingangs aufgez\"ahlten dreizehn Formen.

134. \gesperrt{Relationen zwischen den Formen des Systemes.} Ich will zun\"achst der Uebersichtlichkeit halber wiederum die s\"ammtlichen, bisher gewonnenen Relationen tabellarisch ordnen, und durch einige neue erg\"anzen, die sich direct aus den allgemeinen Formeln \"uber Functionaldeterminanten ergeben.

1) Relationen zwischen den zweiten Ueberschiebungen.
\[
\alpha)\quad
\left.\begin{aligned}
A_{f\vartheta_2} &= A_{f\vartheta_3} = 0\\
A_{\varphi\vartheta_1} &= A_{\varphi\vartheta_3} = 0\\
A_{\psi\vartheta_1} &= A_{\psi\vartheta_2} = 0 % checked against the GDZ scan (27 Oct 2026): the small round dot between \vartheta and the subscript 2 is an ink speck, not a prime; plain A_{\psi\vartheta_2}
\end{aligned}\right\}\ \text{vergl. Nr. 127 (2),}
\]
\[
\beta)\quad
\left.\begin{aligned}
2A_{\vartheta_1\vartheta_1} &= -R_{\varphi\psi} = A_{\varphi\varphi}A_{\psi\psi}-A_{\varphi\psi}^2\\
2A_{\vartheta_2\vartheta_2} &= -R_{f\psi} = A_{ff}A_{\psi\psi}-A_{f\psi}^2\\
2A_{\vartheta_3\vartheta_3} &= -R_{f\varphi} = A_{ff}A_{\varphi\varphi}-A_{f\varphi}^2
\end{aligned}\right\}\ \text{vergl. Nr. 127 (1),}
\]
\[
\gamma)\quad
\left.\begin{aligned}
2A_{\vartheta_1\vartheta_2} &= -A_{\psi\psi}A_{f\varphi}+A_{\psi f}A_{\psi\varphi}\\
2A_{\vartheta_1\vartheta_3} &= -A_{\varphi\varphi}A_{f\psi}+A_{\varphi f}A_{\varphi\psi}
\end{aligned}\right\}\ \text{vergl. Nr. 132 (III).}
\]
\[
2A_{\vartheta_2\vartheta_3} = -A_{ff}A_{\varphi\psi}+A_{f\varphi}A_{f\psi}
\]
2) Relationen zwischen der Combinante und den zweiten Ueberschiebungen:
\[
\alpha)\quad R_{123}=A_{f\vartheta_1}=A_{\varphi\vartheta_2}=A_{\psi\vartheta_3},\ \text{vergl. Nr. 132 (I),}
\]
\[
\beta)\quad 2R_{123}^2=
\begin{vmatrix}
A_{ff} & A_{f\varphi} & A_{f\psi}\\
A_{\varphi f} & A_{\varphi\varphi} & A_{\varphi\psi}\\
A_{\psi f} & A_{\psi\varphi} & A_{\psi\psi}
\end{vmatrix}.
\]
Diese Relation ist eine unmittelbare Folge der allgemeinen Relation (II) Nr. 49. Man erkennt auch, dass die rechten Seiten der Relationen $\beta)$ und $\gamma)$ in (1) nichts anderes als die Minoren dieser Determinante sind. Aus dieser Relation kann man noch zahlreiche andere ableiten, wie sie die Lehrs\"atze \"uber adjungirte Determinanten
%CONT>>


% ===== pdf idx 166, printed 152.0 =====
% !PRINTED 152
\weberpage{152}
%<<CONT
(vergl. Bd. I Nr. 87, 88) liefern; ich will hierauf indess nicht weiter eingehen.

3) Relationen zwischen den Covarianten des Systems.

$\alpha)$ Die Relationen f\"ur die Functionaldeterminanten $(f,\vartheta_i)$, $(\varphi,\vartheta_i)$, $(\psi,\vartheta_i)$, $(\vartheta_i,\vartheta_k)$. Sie ergeben sich sofort aus der allgemeinen Formel Nr. 53 (III). Aus ihr erh\"alt man beispielsweise:
\[
\left.\begin{aligned}
2(f,\vartheta_1) &= -A_{f\varphi}\psi+A_{f\psi}\varphi\\
2(f,\vartheta_2) &= -A_{f\psi}f+A_{ff}\psi\\
2(f,\vartheta_3) &= -A_{ff}\varphi+A_{f\varphi}f
\end{aligned}\right\}\ \text{vergl. auch Nr. 127 (3)}
\]
und ebenso
\[
\left.\begin{aligned}
2(\vartheta_2,\vartheta_1) &= -R_{123}\psi\\
2(\vartheta_1,\vartheta_3) &= -R_{123}\varphi\\
2(\vartheta_3,\vartheta_2) &= -R_{123}f
\end{aligned}\right\}\ \text{vergl. Nr. 133 (IV).}
\]
$\beta)$ Darstellung der Functionaldeterminanten $\vartheta_i$ durch gerade Formen:
\[
\left.\begin{aligned}
-2\vartheta_1^2 &= A_{\varphi\varphi}\psi^2-2A_{\varphi\psi}\varphi\psi+A_{\psi\psi}\varphi^2\\
-2\vartheta_2^2 &= A_{ff}\psi^2-2A_{f\psi}f\psi+A_{\psi\psi}f^2\\
-2\vartheta_3^2 &= A_{ff}\varphi^2-2A_{f\varphi}f\varphi+A_{\varphi\varphi}f^2
\end{aligned}\right\}\ \text{Nr. 48 (I).}
\]
$\gamma)$ Allgemeine Relation f\"ur drei quadratische Formen:
\[
\begin{vmatrix}
A_{ff} & A_{f\varphi} & A_{f\psi} & f\\
A_{f\varphi} & A_{\varphi\varphi} & A_{\varphi\psi} & \varphi\\
A_{f\psi} & A_{\varphi\psi} & A_{\psi\psi} & \psi\\
f & \varphi & \psi & 0
\end{vmatrix}=0
\quad\text{und}\quad
\begin{vmatrix}
A_{\vartheta_1\vartheta_1} & A_{\vartheta_2\vartheta_1} & A_{\vartheta_3\vartheta_1} & \vartheta_1\\
A_{\vartheta_1\vartheta_2} & A_{\vartheta_2\vartheta_2} & A_{\vartheta_3\vartheta_2} & \vartheta_2\\
A_{\vartheta_1\vartheta_3} & A_{\vartheta_2\vartheta_3} & A_{\vartheta_3\vartheta_3} & \vartheta_3\\
\vartheta_1 & \vartheta_2 & \vartheta_3 & 0
\end{vmatrix}=0.
\]
Die beiden letzten Relationen gehen wieder unmittelbar aus der allgemeinen Formel (III) Nr. 50 hervor.

135. \gesperrt{Weitere Relationen zwischen den Formen des Systemes.} Bei der eingreifenden Bedeutung der Theorie quadratischer Formen f\"ur andere Untersuchungen will ich diese Tabelle noch durch einige weitere Relationen erg\"anzen, die wir im Folgenden ableiten wollen. Aus den Identit\"aten
\[
\begin{aligned}
\bar a_0 x_1^2 + 2\bar a_1 x_1 x_2 + a_2 x_2^2 - f &= 0 % [sic] third coefficient printed without a bar, unlike the barred triples in the following three rows
\\
\bar\alpha_0 x_1^2 + 2\bar\alpha_1 x_1 x_2 + \bar\alpha_2 x_2^2 - \varphi &= 0\\
\bar r_0 x_1^2 + 2\bar r_1 x_1 x_2 + \bar r_2 x_2^2 - \psi &= 0\\
y_2^2 x_1^2 - 2y_1y_2x_1x_2 + y_1^2x_2^2 - (xy)^2 &= 0
\end{aligned}
\]
eliminiren wir die vier Gr\"ossen $x_1^2, 2x_1x_2, x_2^2, -1$ und erhalten die nothwendig verschwindende Determinante:
%CONT>>


% ===== pdf idx 167, printed 153.0 =====
% !PRINTED 153
\weberpage{153}

\[
\begin{vmatrix}
\overline{a_0} & \overline{a_1} & \overline{a_2} & f \\
\overline{\alpha_0} & \overline{\alpha_1} & \overline{\alpha_2} & \varphi \\
\overline{r_0} & \overline{r_1} & \overline{r_2} & \psi \\
y_2{}^2 & -y_1y_2 & y_1{}^2 & (xy)^2
\end{vmatrix} = 0.
\]

Ordnen wir sie nach der letzten Colonne, so erhalten wir

\[
R\,(xy)^2 = (\overline{\alpha_0}\,\overline{r_1}\,y_1{}^2)\cdot f - (\overline{a_0}\,\overline{r_1}\,y_1{}^2)\cdot\varphi + (\overline{a_0}\,\overline{\alpha_1}\,y_1{}^2)\cdot\psi\,.
\]

Nun ist aber:

\[
(\overline{\alpha_0}\,\overline{r_1}\,y_1{}^2) =
\begin{vmatrix}
\overline{\alpha_0} & \overline{\alpha_1} & \overline{\alpha_2} \\
\overline{r_0} & \overline{r_1} & \overline{r_2} \\
y_1{}^2 & -y_1y_2 & y_1{}^2
\end{vmatrix}
=
\begin{vmatrix}
\alpha_1{}^2 & \alpha_1\alpha_2 & \alpha_2{}^2 \\
r_1{}^2 & r_1r_2 & r_2{}^2 \\
y_2{}^2 & -y_1y_2 & y_1{}^2
\end{vmatrix}
\]
\[
= (\alpha_1 r_2 - \alpha_2 r_1)(\alpha_1 y_1 + \alpha_2 y_2)(r_1 y_1 + r_2 y_2)
\]
\[
= (\alpha r)\,\alpha_y r_y = \vartheta_1(y)\,.
\]

Demnach besteht folgende Relation:

\weq{I}{R\,(xy)^2 = f(x)\,\vartheta_1(y) + \varphi(x)\,\vartheta_2(y) + \psi(x)\,\vartheta_3(y)}

oder auch, wenn wir $x$ mit $y$ vertauschen,

\weq{II}{R\,(xy)^2 = f(y)\,\vartheta_1(x) + \varphi(y)\,\vartheta_2(x) + \psi(y)\,\vartheta_3(x)\,.}

Beide bilden wieder eine Quelle neuer Relationen. F\"ur $y_i = x_i$ ergiebt sich aus beiden Gleichungen (I) und (II) die Relation:

\weq{III}{f\vartheta_1 + \varphi\vartheta_2 + \psi\vartheta_3 = 0\,.}

Ersetzen wir dagegen in (II) $y_1$ durch $-a_2 (= -b_2)$, $y_2$ durch $a_1 (= -b_1)$ etc., so kommt:

\[
\left.
\begin{aligned}
R f &= A_{ff}\vartheta_1 + A_{f\varphi}\vartheta_2 + A_{f\psi}\vartheta_3 \\
R \varphi &= A_{f\varphi}\vartheta_1 + A_{\varphi\varphi}\vartheta_2 + A_{\varphi\psi}\vartheta_3 \\
R \psi &= A_{f\psi}\vartheta_1 + A_{\varphi\psi}\vartheta_2 + A_{\psi\psi}\vartheta_3
\end{aligned}
\right\}\,. \tag{IV$_\text{a}$}
\]

Ebenso findet man aus der Relation (I), indem man darin die Variable $y$ durch die Coefficienten der Form $\vartheta_i$ ersetzt:

\[
\left.
\begin{aligned}
R \vartheta_1 &= A_{\vartheta_1\vartheta_1} f + A_{\vartheta_1\vartheta_2}\varphi + A_{\vartheta_1\vartheta_3}\psi \\
R \vartheta_2 &= A_{\vartheta_2\vartheta_1} f + A_{\vartheta_2\vartheta_2}\varphi + A_{\vartheta_2\vartheta_3}\psi \\
R \vartheta_3 &= A_{\vartheta_3\vartheta_1} f + A_{\vartheta_3\vartheta_2}\varphi + A_{\vartheta_3\vartheta_3}\psi
\end{aligned}
\right\}\,. \tag{IV$_\text{b}$}
\]

In derselben Weise, wie wir seinerzeit (Nr. 48) die Relation f\"ur $2\vartheta^2(x)$
%CONT>>


% ===== pdf idx 168, printed 154.0 =====
% !PRINTED 154
\weberpage{154}

%<<CONT
ermittelt haben, k\"onnen wir nun auch Relationen f\"ur $2\vartheta^2_{1x}\vartheta^2_{1y}$ etc. berechnen. Es ist n\"amlich:

\[
\vartheta_1(x) = (\varphi,\ \psi) =
\begin{vmatrix}
\overline{\alpha_0} & \overline{\alpha_1} & \overline{\alpha_2} \\
\overline{r_0} & \overline{r_1} & \overline{r_2} \\
x_2{}^2 & -x_1x_2 & x_2{}^2
\end{vmatrix}
\]
\[
2\vartheta_1(y) =
\begin{vmatrix}
\alpha_2 & -2\alpha_1 & \overline{\alpha_0} \\
\overline{r_2} & -2\overline{r_1} & \overline{r_0} \\
y_1{}^2 & 2y_1y_2 & y_2{}^2
\end{vmatrix}\,.
\]

Durch Multiplication beider Determinanten erhalten wir:

\[
2\vartheta_1(x)\vartheta_1(y) =
\begin{vmatrix}
A_{\varphi\varphi} & A_{\varphi\psi} & \varphi(x) \\
A_{\psi\varphi} & A_{\psi\psi} & \psi(x) \\
\varphi(y) & \psi(y) & (xy)^2
\end{vmatrix}\,. \tag{V}
\]

Analoge Relationen hat man f\"ur $2\vartheta_2(x)\vartheta_2(y)$, $2\vartheta_3(x)\vartheta_3(y)$. R\"andern wir ferner die Determinante

\[
R = (\overline{a_0}\,\overline{\alpha_1}\,\overline{r_2})
\]

einmal vertical mit Elementen $0$ und horizontal mit $x_1{}^2$, $-x_1x_2$, $x_2{}^2$, ein zweites Mal, nachdem wir in derselben Determinante die Colonnen gest\"urzt und die zweite Colonne mit $-2$ multiplicirt haben, vertical mit $0$, horizontal mit $y_1{}^2$, $2y_1y_2$, $y_2{}^2$ und bilden das Product der so erhaltenen zwei viergliedrigen Determinanten, so kommt als Resultat:

\[
\begin{vmatrix}
A_{ff} & A_{f\varphi} & A_{f\psi} & f(y) \\
A_{\varphi f} & A_{\varphi\varphi} & A_{\varphi\psi} & \varphi(y) \\
A_{\psi f} & A_{\psi\varphi} & A_{\psi\psi} & \psi(y) \\
f(x) & \varphi(x) & \psi(x) & (xy)^2
\end{vmatrix} = 0\,. \tag{VI}
\]

R\"andern wir endlich die Determinante $R$ durch zwei Saumreihen, n\"amlich einmal vertical durch die Doppelreihe $0\ 0$, und horizontal durch die Zeilen

\[
\begin{aligned}
x_2{}^2 & -x_1x_2 & x_1{}^2 & 0 & 0 \\
y_2{}^2 & -y_1y_2 & y_1{}^2 & 0 & 0\,,
\end{aligned}
\]

ein zweites Mal, nachdem wir die Colonnen gest\"urzt und die zweite Colonne mit $-2$ multiplicirt haben, vertical wieder durch die Doppelreihe $0\ 0$, aber horizontal durch die Zeilen

\[
\begin{aligned}
x_1{}^2 & 2x_1x_2 & x_2{}^2 & 0 & 0 \\
y_1{}^2 & 2y_1y_2 & y_2{}^2 & 0 & 0\,,
\end{aligned}
\]
%CONT>>


% ===== pdf idx 169, printed 155.0 =====
% !PRINTED 155
\weberpage{155}

%<<CONT
so entsteht durch Multiplication der beiden f\"unfgliedrigen Determinanten die letzte Relation:

\[
\begin{vmatrix}
A_{ff} & A_{f\varphi} & A_{f\psi} & f(x) & f(y) \\
A_{\varphi f} & A_{\varphi\varphi} & A_{\varphi\psi} & \varphi(x) & \varphi(y) \\
A_{\psi f} & A_{\psi\varphi} & A_{\psi\psi} & \psi(x) & \psi(y) \\
f(x) & \varphi(x) & \psi(x) & 0 & (xy)^2 \\
f(y) & \varphi(y) & \psi(y) & (xy)^2 & 0
\end{vmatrix} = 0\,. \tag{VII}
\]

Damit m\"ogen die Relationen, welche zwischen den Formen des Systemes stattfinden, ihren Abschluss finden. Wir werden in mannigfachen Untersuchungen von ihnen Gebrauch machen.

136. \textit{Das simultane System von $n$ quadratischen Formen.} Aus dem System der drei Formen $f$, $\varphi$ und $\psi$ l\"asst sich nun leicht das System beliebig vieler quadratischer Formen erschliessen. Bezeichnen wir die Formen mit $f_i$, so besteht dasselbe

1) aus den $n$ Formen $f_1, f_2, f_3 \ldots f_n$ selbst,

2) aus den $\dfrac{n(n-1)}{2}$ Functionaldeterminanten $\vartheta_{ik} = (f_i, f_k)$,

3) aus den $\dfrac{n(n+1)}{2}$ Invarianten $A_{f_i f_k} = (f_i, f_k)^2$,\,*)

4) aus den $\dfrac{n\cdot(n-1)(n-2)}{1\cdot 2\cdot 3}$ Combinanten $R_{ikm} = (f_i, (f_k, f_m))^2 = (f_k, (f_m, f_i))^2 = (f_m, (f_i, f_k))^2$.\,*)

Man k\"onnte noch die Frage aufwerfen, ob nicht die Ueberschiebungen der Functionaldeterminanten zu neuen Formen f\"uhren, n\"amlich zweier solcher Functionaldeterminanten, welche zusammen vier verschiedene Grundformen enthalten, wie

\[
((f_i, f_k), (f_\varrho, f_\sigma))^\lambda\,, \qquad \lambda = 1,\ 2\,. \tag{I}
\]

Doch ist auch deren Reducibilit\"at bereits nachgewiesen. Die erste Ueberschiebung ist n\"amlich reducibel, als Functionaldeterminante von Functionaldeterminanten, die zweite nach Nr. 132 (II). Das System besteht also in der That nur aus den oben angef\"uhrten Formen. Zwischen diesen bestehen nat\"urlich wieder eine grosse Zahl von Relationen, welche den im Vorhergehenden aufgestellten entsprechen,

\noindent\rule{4cm}{0.4pt}

*) Legt man, wie Herr Study es gethan hat, bei der symbolischen Darstellung einer Form $f$ nicht wie wir lineare Factoren zu Grunde, sondern quadratische, so lassen sich alle Invarianten von $f = a_x^{2n}$ darstellen durch simultane Invarianten dieser beiden Formengruppen vom Typus $(ab)^2$ und $(ab)(ac)(bc)$.
%CONT>>


% ===== pdf idx 170, printed 156.0 =====
% !PRINTED 156
\weberpage{156}

%<<CONT
von denen ich hier nur diejenigen anf\"uhren will, die sich auf die unter (I) angegebenen Ueberschiebungen beziehen. Man hat:

\[
1)\quad ((f_i, f_k), (f_\varrho, f_\sigma))^0 = \vartheta_{ik}\cdot\vartheta_{\varrho\sigma}
= -\frac{1}{2}\left\{ A_{i\varrho} f_k f_\sigma - A_{i\sigma} f_k f_\varrho + A_{k\sigma} f_i f_\varrho - A_{k\varrho} f_i f_\sigma \right\}\,, \qquad \text{Nr. 48 (II)}
\]

\[
2)\quad ((f_i, f_k), (f_\varrho, f_\sigma))^1 = (\vartheta_{ik},\ \vartheta_{\varrho\sigma})
= \frac{1}{2}\left\{ R_{i\varrho\sigma} f_k - R_{k\varrho\sigma} f_i \right\} \qquad \text{Nr. 53 (III)}
\]

\[
3)\quad ((f_i, f_k), (f_\varrho, f_\sigma))^2 = (\vartheta_{ik},\ \vartheta_{\varrho\sigma})^2
= \frac{1}{2}\left\{ A_{i\varrho} A_{k\sigma} - A_{i\sigma} A_{k\varrho} \right\}\,. \qquad \text{Nr. 132 (II)}
\]

\[
4)\quad (A_{\varrho\varrho},\ A_{kk},\ A_{\sigma\sigma},\ A_{ii}) = 0\,. \qquad \text{Nr. 50.}
\]

\begin{center}
\S\ 13. \textbf{Specielle quadratische Formen.}
\end{center}

137. \textit{Conjugirte quadratische Formen.} Bei der Cayley'schen Methode der Aufl\"osung einer Gleichung vierten Grades mit Hilfe der Invariantentheorie (vgl. Nr. 175) zeigen sich drei quadratische Formen $\varphi$, $\psi$, $\chi$ (die quadratischen Factoren der Covariante $t$), deren zweite Ueberschiebungen $A_{\varphi\psi}$, $A_{\varphi\chi}$, $A_{\psi\chi}$ identisch verschwinden. Wir wollen allgemein drei solche quadratische Formen $f_0$, $f_1$, $f_2$ conjugirt nennen, sobald ihre simultanen Invarianten $A_{f_0f_1}$, $A_{f_0f_2}$, $A_{f_1f_2}$ identisch null sind. Das volle System dreier conjugirter quadratischer Formen enth\"alt nicht mehr 13, sondern nur 7 selbst\"andige Formen. Die Bedingung $A_{f_if_k} = 0$ bewirkt n\"amlich, dass die Functionaldeterminanten
\[
\vartheta_0 = (f_1, f_2)\,,\qquad \vartheta_1 = (f_0, f_2)\,,\qquad \vartheta_2 = (f_0, f_1)
\]
den Originalformen $f_0$, resp. $f_1$, $f_2$ proportional werden. Denn nach Formel Nr. 134, 3 $(\alpha)$ ist:

\[
2(f_0,\ \vartheta_0) = -A_{f_0f_1}\cdot f_2 + A_{f_0f_2}\cdot f_1\,.
\]

Da aber die beiden Invarianten rechts nach Voraussetzung verschwinden, so ist:

\[
(f_0,\ \vartheta_0) = 0
\]

und ebenso $(f_1,\ \vartheta_1)$, $(f_2,\ \vartheta_2)$. Das Verschwinden dieser Functionaldeterminanten lehrt aber nach Nr. 128, dass die beiden Formen, aus denen sie gebildet wurden, einander proportional sind. Man hat demnach

\[
\begin{aligned}
\vartheta_0 &= c_0 \cdot f_0 \\
\vartheta_1 &= c_1 \cdot f_1 \\
\vartheta_2 &= c_2 \cdot f_2\,.
\end{aligned}
\]

138. \textit{Kanonische Form dreier conjugirter Formen.} Man kann sich fragen: Welchen algebraischen Ausdruck besitzen drei conjugirte
%CONT>>


% ===== pdf idx 171, printed 157.0 =====
% !PRINTED 157
\weberpage{157}
%<<CONT
Formen, wenn man eine derselben in ihrer einfachsten kanonischen
Form $f_0 = 2x_1x_2$ zu Grunde legt? Zur Beantwortung dieser Frage
nehmen wir an $f_1 = \alpha_x^2$ sei eine zu $f_0 = a_x^2$ conjugirte Form, dann
muss, weil nach Vorausetzung % [sic] printed "Vorausetzung"
\[
A_{f_0f_1} = \bar a_0\bar\alpha_2 - 2\bar\alpha_1\bar a_1 + \bar\alpha_0\bar a_2 = 0, \quad \bar a_0 = \bar a_2 = 0, \quad \bar a_1 = 1,
\]
auch $\alpha_1 = 0$ sein, d. h. die betreffende Form $f_1$ muss dem Büschel
angehören:
\[
F = \bar\alpha_0 x_1^2 + \bar\alpha_2 x_2^2.
\]
Sind nun $f_1$ und $f_2$ zwei Formen des Büschels $F$, die nicht nur zu $f_0$,
sondern auch unter einander conjugirt sind, so müssen ausserdem noch
deren Coefficienten der Bedingung genügen:
\[
A_{f_1f_2} = \alpha_0\bar\alpha_2' + \bar\alpha_0'\bar\alpha_2 = 0.
\]
Hieraus ergiebt sich, dass die drei conjugirten Formen in diesem Falle
die Gestalt annehmen:
\[
\begin{aligned}
f_0 &= 2x_1x_2 \\
f_1 &= \bar\alpha_0 x_1^2 + \bar\alpha_2 x_2^2 \\
f_2 &= (\bar\alpha_0 x_1^2 - \bar\alpha_2 x_2^2)\dfrac{\bar\alpha_0'}{\bar\alpha_0}.
\end{aligned}
\]
\textit{Anmerkung.} Drei specielle Formen dieser Art erhalten wir, wenn
wir hierin $\bar\alpha_2 = 1$ setzen und überdies $\bar\alpha_0$ und $\bar\alpha_0'$ so bestimmen, dass
die Discriminanten $A_{f_1f_1}$, $A_{f_2f_2}$ denselben Werth $-2$ annehmen, den
auch $A_{f_0f_0}$ besitzt. Unter diesen Voraussetzungen wird
\[
\bar\alpha_0 = -1, \quad \bar\alpha_0' = i
\]
und somit werden diese drei speciellen Formen:
\[
\begin{aligned}
f_0 &= \phantom{-}2x_1x_2 \\
f_1 &= -(x_1^2 - x_2^2) \\
f_2 &= +i(x_1^2 + x_2^2).
\end{aligned}
\]
Für diese drei Formen berechnen sich die Proportionalitätsconstanten
in den am Schlusse von Nr. 137 erwähnten Gleichungen:
\[
\vartheta_0 = c_0f_0, \quad \vartheta_1 = c_1f_1, \quad \vartheta_2 = c_2f_2
\]
sehr einfach. Denn gerade wegen dieser Beziehungen reducirt sich
die allgemeine Relation Nr. 134 (2, $\alpha$)
\[
R_{012} = A_{f_0\vartheta_0} = A_{f_1\vartheta_1} = A_{f_2\vartheta_2}
\]
auf
\[
R_{012} = c_0A_{f_0f_0} = c_1A_{f_1f_1} = c_2A_{f_2f_2},
\]
%CONT>>


% ===== pdf idx 172, printed 158.0 =====
% !PRINTED 158
\weberpage{158}
%<<CONT
oder, indem wir die Zahlencoefficienten einführen, auf
\[
\left|\begin{array}{ccc} 0 & 1 & 0 \\ -1 & 0 & 1 \\ -i & 0 & i \end{array}\right| = -2c_0 = -2c_1 = -2c_3. % [sic]
\]
Daraus erhalten wir:
\[
c_0 = c_1 = c_2 = \sqrt{-1} = i.
\]
Das Product dieser drei quadratischen Formen ist eine Form sechsten
Grades
\[
t = \mathrm{Const.}\ x_1x_2(x_1^4 - x_2^4),
\]
welche Klein (vgl. Vorles. über das Ikosaeder § 10) Octaeder nennt.
Man kann nämlich die Gauss'sche Zahlenebene so auf eine im Nullpunkt
berührende Kugel projiciren, dass die sechs Wurzeln der Gleichung
$t = 0$ gerade die Eckpunkte eines in die Kugel einbeschriebenen
Octaeders werden. Wir werden auf das Octaeder noch später zu
sprechen kommen (vgl. § 19).

139. \textit{Die quadratischen Formen des Würfels.} Wir hatten uns in
Nr. 138 die Aufgabe vorgelegt, quadratische Formen aufzustellen,
deren zweite Ueberschiebungen $A_{f_if_k}$ $(i \gtrless k)$ verschwinden. Die Frage
kann man verallgemeinern: Lässt sich eine Gruppe von quadratischen
Formen, die keinen gemeinschaftlichen Factor besitzen, so bestimmen,
dass ihre sämmtlichen zweiten Ueberschiebungen $A_{f_if_k}$ $(i \gtrless k)$ einen
festen Werth $\varrho$ besitzen?

Zur Beantwortung wählen wir wiederum als eine der quadratischen
 Formen die Form
\[
f_0 = 2x_1x_2.
\]
Wir wissen, dass ihre Discriminante den Werth $-2$ besitzt, und wollen
annehmen, dass auch die übrigen Formen $f_i$ $(i = 1, 2 \dots n)$ eine
Normalform besitzen mögen, so dass ihre Discriminante
\[
A_{f_if_i} = -2
\]
sei. Dann lautet die Frage: Welches sind unter diesen Voraussetzungen
 die Formen $f_i$, deren zweite Ueberschiebungen $A_{f_if_k}$ $(i \gtrless k)$
sowohl unter einander als mit $f_0$ den festen Werth $\varrho$ besitzen? Es
mag von vornherein bemerkt sein, dass dieser Werth $\varrho$ von $\pm 2$ verschieden
 sein muss. Denn da die gesuchten Formen keinen gemeinschaftlichen
 Factor besitzen sollen, so müssen die Resultanten
\[
R_{f_if_k} \lessgtr 0
\]
sein. Nun ist aber nach Nr. 134 (1, $\beta$)
%CONT>>


% ===== pdf idx 173, printed 159.0 =====
% !PRINTED 159
\weberpage{159}
%<<CONT
\[
R_{f_if_k} = \left|\begin{array}{cc} A_{f_if_k} & A_{f_if_i} \\ A_{f_kf_k} & A_{f_if_k} \end{array}\right| = \left|\begin{array}{cc} A_{f_if_k} & -2 \\ -2 & A_{f_if_k} \end{array}\right| = A_{f_if_k}^2 - 4 \gtrless 0.
\]
d. h. $A_{f_if_k}$ muss von $\pm 2$ verschieden sein.

140. Sei nun $f_1 = x^2 + 2bx + c$, wobei wir noch der Einfachheit
 halber den Coefficienten von $x^2$ gleich 1 annahmen, eine erste
Form von den verlangten Eigenschaften, und
\[
f_2 = \alpha x^2 + 2\beta x + \gamma
\]
eine zweite, dann müssen nach unseren Voraussetzungen die Beziehungen
 bestehen:
\weq{I}{
\left.\begin{aligned}
A_{f_1f_1} &= 2(c - b^2) = -2 \\
A_{f_0f_1} &= -2b = \varrho \\
A_{f_2f_2} &= 2(\alpha\gamma - \beta^2) = -2 \\
A_{f_0f_2} &= -2\beta = \varrho \\
A_{f_1f_2} &= \gamma + c\alpha - 2b\beta = \varrho
\end{aligned}\right\} .
}
Dies sind gerade fünf Gleichungen mit fünf Unbekannten. Aus den
ersten beiden bestimmen sich die Coefficienten von $f_1$ linear; aus den
drei letzten dagegen erhält man eine quadratische Gleichung für $\alpha$
oder $\gamma$, und erkennt somit, dass es ausser $f_0 f_1 f_2$ noch eine weitere
Form $f_3$ giebt, so dass die fünf Invarianten
\[
A_{f_0f_1}, \quad A_{f_0f_2}, \quad A_{f_0f_3}, \quad A_{f_1f_2}, \quad A_{f_1f_3}
\]
den Werth $\varrho$ besitzen. Das gestellte Problem verlangt aber, dass
auch $A_{f_2f_3} = \varrho$ sei. Da jedoch die Coefficienten von $f_2$ und $f_3$ durch
die obigen fünf Gleichungen schon völlig bestimmt sind, so kann
diese Forderung nur mehr eine Bedingung für $\varrho$ sein. Nehmen wir
also an, $\varrho$ sei eine positive Grösse und fragen uns, welchen numerischen
 Werth muss sie besitzen, damit allgemein
\[
A_{f_if_k} = \varrho
\]
ist. Die Beantwortung der Frage liefert uns die Relation Nr. 136 (4),
welche zwischen den zweiten Ueberschiebungen von vier Formen bestehen
 muss, und die wir allgemein in Nr. 50 abgeleitet haben. Dieselbe
 war:
\[
\left|\begin{array}{cccc}
(f,f)^2, & (f,\varphi)^2, & (f,\psi)^2, & (f,\chi)^2 \\
(\varphi,f)^2, & (\varphi,\varphi)^2, & (\varphi,\psi)^2, & (\varphi,\chi)^2 \\
(\psi,f)^2, & (\psi,\varphi)^2, & (\psi,\psi)^2, & (\psi,\chi)^2 \\
(\chi,f)^2, & (\chi,\varphi)^2, & (\chi,\psi)^2, & (\chi,\chi)^2
\end{array}\right| = 0.
\]


% ===== pdf idx 174, printed 160.0 =====
% !PRINTED 160
\weberpage{160}
Setzen wir in ihr die Werthe $+\varrho$ und $-2$ der zweiten Ueberschiebungen
 ein, so kommt:
\[
\left|\begin{array}{cccc}
-2, & \varrho, & \varrho, & \varrho \\
\varrho, & -2, & \varrho, & \varrho \\
\varrho, & \varrho, & -2, & \varrho \\
\varrho, & \varrho, & \varrho, & -2
\end{array}\right| = 0,
\]
oder, wenn man $-\dfrac{2}{\varrho} = \sigma$ setzt und geeignet reducirt:
\[
(1 - \sigma)^3 \left|\begin{array}{cccc}
\sigma & 1 & 1 & 1 \\
1 & -1 & 0 & 0 \\
1 & 0 & -1 & -1 \\
1 & 0 & 0 & -1
\end{array}\right| = 0,
\]
oder endlich: $(1 - \sigma)^3(\sigma + 3) = 0.$

141. Der cubische Factor $(1 - \sigma)^3 = 0$ liefert $\varrho = -2$ als Lösung;
diese haben wir gerade ausgeschlossen. Es bleibt als einzige Lösung
\[
\sigma = -3
\]
also
\[
\varrho = +\frac{2}{3}.
\]
Unter Benutzung dieses Werthes gehen die fünf Coefficientenbedingungen
(I) über in:
\[
\begin{aligned}
A_{f_1f_1} &= \phantom{-}c - b^2 = -1 \\
A_{f_0f_1} &= -b = \phantom{-}\dfrac{1}{3} \\
A_{f_2f_2} &= \phantom{-}\alpha\gamma - \beta^2 = -1 \\
A_{f_0f_2} &= -\beta = \phantom{-}\dfrac{1}{3} \\
A_{f_1f_2} &= \gamma + c\alpha - 2b\beta = \phantom{-}\dfrac{2}{3}.
\end{aligned}
\]
Daraus folgt zunächst $b = -\dfrac{1}{3}$, $c = -\dfrac{8}{9}$, $\beta = -\dfrac{1}{3}$. Demnach
erhält man für $\alpha$ und $\gamma$:
\weq{II}{
\left.\begin{aligned}
\gamma - \dfrac{8}{9}\alpha &= +\dfrac{8}{9} \\
\alpha\gamma &= -\dfrac{8}{9}
\end{aligned}\right\} .
}
Denken wir uns in $f_2$ die Constante $\gamma$ in der Form $-\dfrac{8}{9}\gamma'$, dann
gehen die Gleichungen (II) in die symmetrisch gebauten:
\[
\gamma' + \alpha = -1
\]
\[
\alpha\gamma' = \phantom{-}1
\]


% ===== pdf idx 175, printed 161.0 =====
